\section*{Methodologie}
\subsection*{Vereisten}
Voor dit onderzoek zijn de volgende zaken het doel: 

1. **Een library ontwikkelen voor algemene bruikbaarheid**: Het doel is om een library te maken die in elke applicatie kan worden gebruikt. Deze library moet het proces van permissies aanvragen, initialisatie, gebruik en uitbreiding zo soepel mogelijk laten verlopen. Dit betekent dat de library flexibel, modulair en eenvoudig te integreren moet zijn in verschillende soorten applicaties.

2. **Een testapplicatie ontwikkelen voor toekomstige onderzoekers**: Naast de library is het doel om een testapplicatie te schrijven. Deze testapplicatie zal dienen als een referentie-implementatie en als hulpmiddel voor toekomstige onderzoekers. Hierdoor hoeven zij niet telkens zelf een applicatie te ontwikkelen, maar kunnen ze eenvoudig hun eigen algoritmes testen in een bestaande, goed gedefinieerde omgeving.

Door deze doelen te realiseren, wordt niet alleen de bruikbaarheid van de ontwikkelde oplossingen vergroot, maar wordt ook de drempel voor toekomstig onderzoek verlaagd.
\subsection*{Veronderstellingen}
\subsection*{Flow diagram}
